% =====================================================================
%  Fiche pédagogique — Physique, chapitre 5 : L’énergie thermique, le transfert thermique (4 h, deux séances)
% =====================================================================
\newcommand{\Matiere}{Physique}
\newcommand{\Partie}{Partie 1 — Travail mécanique et énergie}
\newcommand{\NumChap}{5}
\newcommand{\TitreChap}{L’énergie thermique :\\ le transfert thermique}
\newcommand{\TitreCourt}{L’énergie thermique : le transfert thermique}
\newcommand{\Niveau}{1\textsuperscript{re} année du baccalauréat}
\newcommand{\Duree}{4 h}
\newcommand{\Banniere}{opener_phys5.jpg}
\newcommand{\BanniereDecalage}{1.3cm}
\newcommand{\dC}[1]{\qty{#1}{\degreeCelsius}}
\newcommand{\cu}{\unit{J.kg^{-1}.\degreeCelsius^{-1}}}
\newcommand{\Em}{E_{\mathrm{m}}}

\begin{document}
\bandeau

\noindent
\begin{minipage}[t]{0.235\linewidth}\infocarte{fi_toque}{cPurple}{Niveau}{\Niveau\\ Sciences expérimentales (option française)}\end{minipage}\hfill
\begin{minipage}[t]{0.235\linewidth}\infocarte{fi_horloge}{cAmber}{Durée}{4 heures — deux séances\\ 2 h + 2 h}\end{minipage}\hfill
\begin{minipage}[t]{0.235\linewidth}\infocarte{fi_prof}{cCyan}{Professeur}{\rule{0pt}{3.6mm}{\color{cInk2!60}\dotfill}\\ Classe : {\color{cInk2!60}\dotfill}}\end{minipage}\hfill
\begin{minipage}[t]{0.235\linewidth}\infocarte{fi_chemin}{cInk}{Programme}{Physique · Partie 1 · § 4\\ L’énergie thermique}\end{minipage}

% =====================================================================
\bloc{1}{fi_cible}{Cadre pédagogique de la séquence}

\begin{xltabular}{\linewidth}{@{}C{3.5cm}I Y@{}}
\ligne{fi_cible}{cRed}{Compétence\\ spécifique}{%
\textbf{Travail mécanique et énergie :} \textbf{interpréter les transferts d’énergie} et les phénomènes de \textbf{conservation et de dissipation} dans diverses situations de la vie quotidienne ; \textbf{résoudre une situation-problème} relative à la conservation et à la dissipation de l’énergie dans un système mécanique, expérimentalement ou par une étude analytique {\color{cInk2}(compétence spécifique de la physique, 1\textsuperscript{re} année du baccalauréat)}.
\pt{\textbf{Contribution de ce chapitre :} la \textbf{chaleur}, forme de \textbf{transfert d’énergie} ; quantité de chaleur, chaleur massique, capacité thermique ; \textbf{équilibre thermique} et \textbf{équation calorimétrique} ; chaleur latente de changement d’état ; le \textbf{rayonnement}, autre forme de transfert d’énergie. Situations-problèmes : chauffe-eau, calorimètre, fusion de la glace, four solaire, bouteille isolante.}}
\ligne{fi_reseau}{cTeal}{Compétences\\ transversales}{%
\begin{minipage}[t]{0.493\linewidth}\RaggedRight
\pt{\textbf{Démarche scientifique} — maîtriser les éléments de la démarche scientifique dans toutes ses dimensions.\hfill\break{\color{cTeal}\itshape Ici :} expériences de chauffage (une grandeur varie à la fois) ; mesures au calorimètre ; comparaison avec les valeurs admises et discussion des écarts.}
\pt{\textbf{Communication} — lire, écrire, écouter, dialoguer ; maîtriser la langue écrite et le langage mathématique ; utiliser les symboles ; respecter l’éthique du dialogue.\hfill\break{\color{cTeal}\itshape Ici :} bilans calorimétriques écrits avec rigueur (signes, unités) ; lecture d’une courbe d’échauffement.}
\pt{\textbf{Technologies de l’information} — utiliser les TIC pour s’informer, traiter des données, éprouver des hypothèses, construire des modèles.\hfill\break{\color{cTeal}\itshape Ici :} animations (chauffage, calorimètre, courbe de la glace, rayonnement) ; thermomètres numériques.}
\end{minipage}\hfill
\begin{minipage}[t]{0.483\linewidth}\RaggedRight
\pt{\textbf{Méthodologie} — acquérir une méthodologie de recherche, de travail et d’autoformation ; réaliser des tâches complémentaires.\hfill\break{\color{cTeal}\itshape Ici :} méthode en cinq étapes pour résoudre un problème de calorimétrie (« finale moins initiale », vérification de $\theta_e$).}
\pt{\textbf{Culture intégrée} — s’intégrer dans son environnement, citoyenneté active, positions positives face aux questions d’environnement, de santé, de prévention et de consommation.\hfill\break{\color{cTeal}\itshape Ici :} économies d’énergie (chauffe-eau, isolation, bouteille isolante) ; énergie solaire (chauffe-eau solaires, fours solaires) ; climat côtier ; sécurité avec l’eau bouillante.}
\end{minipage}}
\ligne{fi_drapeau}{cBlue}{Objectifs\\ d’apprentissage}{%
{\color{cInk2}\itshape À la fin de la séquence, l’apprenant(e) est capable de :}\par\vspace{.6mm}
\begin{minipage}[t]{0.493\linewidth}\RaggedRight
\pt{\textbf{O1} — savoir que la \textbf{chaleur} est une forme de transfert d’énergie, du corps chaud vers le corps froid ;}
\pt{\textbf{O2} — connaître et utiliser l’expression $Q = m\cdot c\cdot\Delta\theta$, son unité et la convention de signe ;}
\pt{\textbf{O3} — définir la \textbf{chaleur massique} d’un corps et la \textbf{capacité thermique} d’un calorimètre, avec leurs unités ;}
\pt{\textbf{O4} — reconnaître l’\textbf{équilibre thermique} et appliquer l’\textbf{équation calorimétrique} $\sum Q_i = 0$ ;}
\end{minipage}\hfill
\begin{minipage}[t]{0.483\linewidth}\RaggedRight
\pt{\textbf{O5} — connaître la \textbf{chaleur latente} de changement d’état et son unité ; utiliser $Q = \pm m\cdot L$ ;}
\pt{\textbf{O6} — \textbf{déterminer expérimentalement} la capacité thermique d’un calorimètre, la chaleur massique d’un métal et la chaleur latente de fusion de la glace ;}
\pt{\textbf{O7} — reconnaître le \textbf{rayonnement} comme une forme de transfert d’énergie et le distinguer de la conduction et de la convection.}
\end{minipage}
\par\vspace{1mm}{\footnotesize\color{cInk2}\textit{Connaissances et habiletés du programme :} savoir que la chaleur est une forme de transfert d’énergie ; connaître l’expression de la quantité de chaleur $Q = m\cdot c\cdot\Delta\theta$ et son unité ; reconnaître l’équilibre thermique et appliquer la relation qui l’exprime ; connaître la chaleur massique d’un métal et son unité ; connaître la chaleur latente de changement d’état et son unité ; déterminer la capacité thermique, la chaleur massique et la chaleur latente ; reconnaître le rayonnement comme une forme de transfert d’énergie.}}
\ligne{fi_marches}{cGreen}{Prérequis}{%
\begin{minipage}[t]{0.493\linewidth}\RaggedRight
{\sffamily\bfseries\footnotesize\color{cGreen}Acquis du collège, du tronc commun et du chapitre 4}
\pt{température en degrés Celsius et en kelvins ($T = \theta + 273{,}15$) ; mesure d’une température ;}
\pt{états de la matière et changements d’état (fusion, solidification, vaporisation, liquéfaction) ;}
\pt{masse volumique $\rho = m/V$ ; énergie et puissance $E = P\cdot\Delta t$ ;}
\pt{dissipation de l’énergie mécanique en chaleur : $\Delta\Em = -Q$ (chapitre 4).}
\end{minipage}\hfill
\begin{minipage}[t]{0.483\linewidth}\RaggedRight
{\sffamily\bfseries\footnotesize\color{cGreen}Outils à réactiver}
\pt{proportionnalité ; résolution d’une équation du premier degré ;}
\pt{lecture d’une courbe $\theta = f(t)$ (paliers) ;}
\pt{conversions (\unit{g} $\to$ \unit{kg}, \unit{L} $\to$ \unit{kg}, \unit{min} $\to$ \unit{s}, \unit{kJ}, \unit{MJ}) ; écriture scientifique et chiffres significatifs.}
\end{minipage}}
\ligne{fi_boite}{cOrange}{Matériel\\ didactique}{%
\begin{minipage}[t]{0.493\linewidth}\RaggedRight
{\sffamily\bfseries\footnotesize\color{cOrange}Expériences de cours}
\pt{deux réchauds ou plaques électriques identiques, béchers, thermomètres numériques, chronomètre, balance ; eau, huile végétale ;}
\pt{calorimètre de Berthelot avec agitateur et thermomètre ; eau tiède ; cylindre d’aluminium, bain-marie (eau bouillante), pince ;}
\pt{glace fondante et glaçons, papier absorbant ;}
\pt{cloche à vide, platine et pompe ; thermomètre à réservoir noirci ; lampe infrarouge ; plaques métalliques peintes en noir mat et en blanc brillant ; bouteille isolante.}
\end{minipage}\hfill
\begin{minipage}[t]{0.483\linewidth}\RaggedRight
{\sffamily\bfseries\footnotesize\color{cOrange}Supports}
\pt{vidéoprojecteur et ordinateur : \textbf{présentation animée} du chapitre (38 diapositives) ;}
\pt{\textbf{cahier} de l’apprenant : fiche d’activités collée, réponses aux questions, « À retenir », applications et exercices ;}
\pt{\textbf{fiche d’activités} imprimée (deux pages : activités 1 et 2, puis 3 et 4, à coller dans le cahier) ;}
\pt{tableau, feutres ou craies de couleur, calculatrices.}
\end{minipage}}
\ligne{fi_livre}{cInk2}{Références}{%
\pt{Ministère de l’Éducation nationale — \textit{Orientations pédagogiques générales et programmes spécifiques à l’enseignement de la physique et de la chimie au cycle secondaire qualifiant}, 2007 : programme de la 1\textsuperscript{re} année du baccalauréat — physique, § 4 « L’énergie thermique : le transfert thermique » (propre aux filières sciences expérimentales et sciences et technologies) : chaleur massique d’un corps pur ; quantité de chaleur $Q = m\cdot c\cdot\Delta\theta$ et son signe conventionnel ; équilibre thermique, équation calorimétrique ; chaleur latente de changement d’état d’un corps pur ; une autre forme de transfert d’énergie : le rayonnement.}
\pt{Activités proposées : expériences simples pour mettre en évidence les grandeurs dont dépend la quantité de chaleur ; étude expérimentale quantitative du transfert thermique entre deux corps pour déterminer la capacité thermique d’un calorimètre, la chaleur massique d’un métal et la chaleur latente de changement d’état d’un solide ; expériences, exemples de l’environnement de l’apprenant, documents ou logiciels pour montrer que le rayonnement est une autre forme de transfert d’énergie.}
\pt{Orientations : montrer par des expériences simples que la quantité de chaleur dépend de la masse, de la nature de la substance et de la variation de température, puis donner son expression ; définir la chaleur massique d’un corps pur, la capacité thermique d’un calorimètre et la chaleur latente de changement d’état d’un solide ; indiquer que le bilan calorimétrique ne dépend que des états initial et final. Supports : présentation animée et fiche d’activités ; manuels scolaires agréés de la 1\textsuperscript{re} année du baccalauréat.}}
\end{xltabular}

% =====================================================================
\needspace{7cm}
\bloc{2}{fi_chemin}{Déroulement de la séquence (deux séances)}

\seance{1}{2 h}{quantité de chaleur et chaleur massique ; équilibre thermique et calorimétrie}
\frise[120]{cCyan/10/Mise en route, cTeal/53/Partie 1 · Chaleur massique, cBlue/52/Partie 2 · Calorimétrie, cGreen/5/Bilan}
\seance{2}{2 h}{chaleur latente ; rayonnement ; évaluation}
\frise[120]{cCyan/8/Rappels, cPurple/52/Partie 3 · Changement d’état, cAmber/38/Partie 4 · Rayonnement, cOrange/22/Évaluation et bilan}
\par\vspace{4mm}

\begin{xltabular}{\linewidth}{@{}L{3.75cm}I >{\small\RaggedRight\arraybackslash}XI >{\small\RaggedRight\arraybackslash}XI L{4.6cm}I C{1.45cm}@{}}
\enteteDeroulement
\endfirsthead
\enteteDeroulement
\endhead
% =====================================================================
\partie{cInk}{Séance 1 · Chaleur massique et calorimétrie (2 h)}{0 h 00 $\to$ 2 h 00}
\etape{cCyan}{\axe{cInk2}{ic_rappel}{Mise en route}{Accueil, plan et prérequis}{\diapo{1} \support{Fiche}}}
{\pt{Distribue la fiche d’activités (deux pages), que chacun colle dans son cahier ; projette la diapositive de titre et annonce le plan des deux séances.}
 \pt{Questions flash : « \dC{25} en kelvins ? », « masse de \qty{2}{L} d’eau ? », « énergie fournie par un appareil de \qty{100}{W} en \qty{1}{min} ? », « nom du passage de l’état solide à l’état liquide ? ».}}
{\pt{Répond oralement : \qty{298}{K} ; \qty{2}{kg} ; \qty{6000}{J} ; la fusion.}
 \pt{Colle la fiche d’activités dans le cahier.}}
{\diag Questions flash.\par
 \textit{Indicateurs :} conversions ; $E = P\cdot\Delta t$ ; vocabulaire des changements d’état.\par\vspace{.8mm}
 \regul Rappel au tableau des changements d’état et de $\rho = m/V$.}
{10}{0 h 10}
% ---------------------------------------------------------------------
\partie{cTeal}{Partie 1 · Quantité de chaleur et chaleur massique}{53 min · objectifs O1 à O3}
\etape{cTeal}{\axe{cAmber}{ic_situation}{Situation de départ}{Grande casserole et petite tasse}{\diapo{2}}}
{\pt{Projette l’animation (diapo 2) : deux masses d’eau chauffées sur des réchauds identiques.}
 \pt{Pose le problème : de quoi dépend la chaleur reçue par un corps qui s’échauffe ?}}
{\pt{Propose : « la masse », « l’écart de température », « la matière ».}}
{\diag Hypothèses des apprenants.}
{5}{0 h 15}
\etape{cTeal}{\axe{cTeal}{ic_activite}{Activité 1}{De quoi dépend la chaleur reçue ? (docs 1 à 3)}{\diapo{3--6} \support{Fiche · Act. 1} \support{Cahier}}}
{\pt{Réalise devant la classe le chauffage de \qty{100}{g} et \qty{200}{g} d’eau sur deux réchauds identiques (doc 1, diapo 3), puis compare l’eau et l’huile (doc 3, diapo 4).}
 \pt{Fait remarquer que l’énergie reçue est proportionnelle à la durée de chauffage ; fait traiter les questions 1 à 4 en binômes ; corrige au clic (diapos 5 et 6).}}
{\pt{Travaille sur la fiche d’activités collée dans le cahier ; rédige ses réponses dans le cahier.}
 \pt{Établit $Q \propto m$ et $Q \propto \Delta\theta$ ; constate que l’huile s’échauffe deux fois plus que l’eau ; écrit $Q = m\cdot c\cdot(\theta_f - \theta_i)$, en joules ; définit $c$, $C = m\cdot c$ et le signe de $Q$.}}
{\form Observation des binômes, correction collective.\par
 \textit{Indicateurs :} raisonnement de proportionnalité ; unités de $c$.\par\vspace{.8mm}
 \regul Erreur fréquente : calculer $\theta_i - \theta_f$ (rappeler « finale moins initiale »).}
{30}{0 h 45}
\etape{cTeal}{\axe{cRed}{ic_retenir}{À retenir}{$Q = m\cdot c\cdot\Delta\theta$ ; capacité thermique}{\diapo{7--8} \support{Cahier}}}
{\pt{Fait formuler la définition de la chaleur et l’expression de $Q$ ; projette les diapos 7 et 8 (tableau des chaleurs massiques) et laisse le temps de recopier.}}
{\pt{Recopie l’« À retenir » dans son cahier.}}
{\form Question flash : « quel corps du tableau s’échauffe le plus vite pour la même énergie ? » $\to$ le plomb.}
{10}{0 h 55}
\etape{cTeal}{\axe{cPurple}{ic_application}{Application}{Un chauffe-eau électrique}{\diapo{9} \support{Cahier}}}
{\pt{Projette l’énoncé (diapo 9) ; travail individuel dans le cahier ; correction au clic.}}
{\pt{Calcule $m = \qty{100}{kg}$, $C \approx \qty{4,19e5}{J.\degreeCelsius^{-1}}$, $Q \approx \qty{2,09e7}{J}$ ; justifie $Q > 0$.}}
{\somm Exercice individuel.\par
 \textit{Indicateurs :} $m = \rho\cdot V$ ; écriture scientifique.}
{8}{1 h 03}
% ---------------------------------------------------------------------
\partie{cBlue}{Partie 2 · Équilibre thermique et calorimétrie}{52 min · objectifs O4 et O6}
\etape{cBlue}{\axe{cAmber}{ic_situation}{Situation de départ}{Un métal chaud dans l’eau froide}{\diapo{10}}}
{\pt{Projette l’animation (diapo 10) : les températures du métal et de l’eau tendent vers une même valeur.}
 \pt{Pose le problème : comment mesurer la capacité thermique d’un récipient et la chaleur massique d’un métal ?}}
{\pt{Propose : « utiliser un récipient isolé », « mesurer la température finale ».}}
{\diag Idée d’équilibre thermique.}
{4}{1 h 07}
\etape{cBlue}{\axe{cTeal}{ic_activite}{Activité 2}{Mesurer avec un calorimètre (doc 4)}{\diapo{11--14} \support{Fiche · Act. 2} \support{Cahier}}}
{\pt{Présente le calorimètre (diapo 11) ; réalise le protocole 1 (eau froide + eau tiède), puis le protocole 2 (cylindre d’aluminium sorti de l’eau bouillante, avec une pince) ; fait relever les températures d’équilibre.}
 \pt{Fait traiter les questions 1 à 4 ; corrige au clic (diapos 13 et 14).}}
{\pt{Écrit $Q_1 = (\mu + m_1\cdot c_e)(\theta_e - \theta_1)$ et $Q_2 = m_2\cdot c_e\cdot(\theta_e - \theta_2)$ ; trouve $\mu \approx \qty{109}{J.\degreeCelsius^{-1}}$.}
 \pt{Trouve $c_{\mathrm{Al}} \approx \qty{892}{J.kg^{-1}.\degreeCelsius^{-1}}$ et le compare à 897.}}
{\form Observation des binômes ; correction collective.\par
 \textit{Indicateurs :} le calorimètre est compté dans le bilan ; signes de $Q_1$ et $Q_2$.\par\vspace{.8mm}
 \regul Erreur fréquente : oublier $\mu$, ou écrire $Q_{\text{reçue}} = Q_{\text{cédée}}$ sans valeurs absolues.}
{30}{1 h 37}
\etape{cBlue}{\axe{cRed}{ic_retenir}{À retenir et méthode}{$\sum Q_i = 0$ ; méthode en cinq étapes}{\diapo{15--16} \support{Cahier}}}
{\pt{Fait énoncer l’équation calorimétrique ; projette les diapos 15 et 16 ; laisse le temps de recopier.}}
{\pt{Recopie l’« À retenir » et la méthode dans son cahier.}}
{\form Question flash : « $\theta_e$ peut-elle dépasser la plus haute température initiale ? » $\to$ non.}
{10}{1 h 47}
\etape{cBlue}{\axe{cPurple}{ic_application}{Application}{Un échantillon de cuivre}{\diapo{17} \support{Cahier}}}
{\pt{Projette l’énoncé (diapo 17) ; travail individuel en suivant la méthode ; correction au clic.}}
{\pt{Écrit les deux chaleurs en fonction de $\theta_e$ ; trouve $\theta_e \approx \dC{20,2}$.}}
{\somm Exercice individuel.\par
 \textit{Indicateur :} vérification $\dC{16,0} < \theta_e < \dC{95,0}$.}
{8}{1 h 55}
\etape{cGreen}{\axe{cGreen}{ic_reponses}{Bilan de la séance 1}{Synthèse et travail à faire}{\diapo{15--16}}}
{\pt{Fait résumer oralement la séance ; donne le travail à faire (exercice 1, énoncé projeté, diapo 32).}}
{\pt{Résume : $Q = m\cdot c\cdot\Delta\theta$, $\sum Q_i = 0$ ; recopie l’énoncé de l’exercice 1.}}
{\form Synthèse orale.}
{5}{2 h 00}
% =====================================================================
\partie{cInk}{Séance 2 · Changement d’état, rayonnement et évaluation (2 h)}{0 h 00 $\to$ 2 h 00}
\etape{cCyan}{\axe{cInk2}{ic_rappel}{Rappels}{Correction de l’exercice 1}{\diapo{32--33}}}
{\pt{Corrige l’exercice 1 au clic (diapo 33) ; questions flash sur $Q = m\cdot c\cdot\Delta\theta$ et $\sum Q_i = 0$.}}
{\pt{Vérifie : $\mu \approx \qty{105}{J.\degreeCelsius^{-1}}$ ; $c_{\mathrm{Pb}} \approx \qty{135}{J.kg^{-1}.\degreeCelsius^{-1}}$.}}
{\diag Vérification du travail personnel.\par\vspace{.8mm}
 \regul Reprendre la méthode si besoin.}
{8}{0 h 08}
% ---------------------------------------------------------------------
\partie{cPurple}{Partie 3 · Chaleur latente de changement d’état}{52 min · objectifs O5 et O6}
\etape{cPurple}{\axe{cAmber}{ic_situation}{Situation de départ}{La glace qui fond}{\diapo{18}}}
{\pt{Projette la courbe animée (diapo 18) : la température reste à \dC{0} pendant la fusion.}
 \pt{Pose le problème : comment calculer l’énergie reçue si $\Delta\theta = 0$ ?}}
{\pt{Constate le palier ; propose : « l’énergie sert à faire fondre la glace ».}}
{\diag Représentations sur le palier.}
{4}{0 h 12}
\etape{cPurple}{\axe{cTeal}{ic_activite}{Activité 3}{L’énergie d’un changement d’état (docs 5 et 6)}{\diapo{19--22} \support{Fiche · Act. 3} \support{Cahier}}}
{\pt{Exploite la courbe (doc 5, diapo 19) ; introduit $Q = m\cdot L$.}
 \pt{Réalise la mesure de $L_f$ : glaçon de glace fondante, essuyé, introduit dans le calorimètre (doc 6, diapo 20) ; fait traiter les questions 1 à 5 ; corrige au clic (diapos 21 et 22).}}
{\pt{Décrit les trois parties de la courbe ; définit $L$ (\unit{J.kg^{-1}}) ; écrit $Q_g = m_2\cdot L_f + m_2\cdot c_e\cdot(\theta_e - 0)$ et $Q_{ce} = (\mu + m_1\cdot c_e)(\theta_e - \theta_1)$.}
 \pt{Calcule $L_f \approx \qty{3,35e5}{J.kg^{-1}}$.}}
{\form Correction collective.\par
 \textit{Indicateurs :} deux étapes pour le glaçon ; signe de $Q = \pm m\cdot L$.\par\vspace{.8mm}
 \regul Erreur fréquente : oublier l’échauffement de l’eau de fusion.}
{30}{0 h 42}
\etape{cPurple}{\axe{cRed}{ic_retenir}{À retenir}{$Q = \pm m\cdot L$}{\diapo{23} \support{Cahier}}}
{\pt{Projette la diapo 23 (courbe légendée, tableau des chaleurs latentes) ; laisse le temps de recopier.}}
{\pt{Recopie l’« À retenir » dans son cahier.}}
{\form Question flash : « signe de $Q$ pour la solidification de l’eau ? » $\to$ négatif.}
{8}{0 h 50}
\etape{cPurple}{\axe{cPurple}{ic_application}{Application}{De la glace à l’eau tiède}{\diapo{24} \support{Cahier}}}
{\pt{Projette l’énoncé (diapo 24) ; travail individuel ; correction au clic.}}
{\pt{Identifie les trois étapes ; calcule $Q_1 \approx \qty{1,57e4}{J}$, $Q_2 = \qty{1,67e5}{J}$, $Q_3 \approx \qty{4,19e4}{J}$, $Q_{\mathrm{tot}} \approx \qty{2,25e5}{J}$.}}
{\somm Exercice individuel.\par
 \textit{Indicateurs :} bonne chaleur massique pour chaque étape ($c_g$ ou $c_e$).}
{10}{1 h 00}
% ---------------------------------------------------------------------
\partie{cAmber}{Partie 4 · Une autre forme de transfert d’énergie : le rayonnement}{38 min · objectif O7}
\etape{cAmber}{\axe{cAmber}{ic_situation}{Situation de départ}{La chaleur du Soleil}{\diapo{25}}}
{\pt{Projette l’animation (diapo 25) : l’énergie du Soleil traverse le vide spatial.}
 \pt{Pose le problème : comment l’énergie se propage-t-elle sans matière ?}}
{\pt{Propose : « par la lumière », « par des rayons ».}}
{\diag Représentations initiales.}
{4}{1 h 04}
\etape{cAmber}{\axe{cTeal}{ic_activite}{Activité 4}{Un transfert d’énergie sans matière (docs 7 et 8)}{\diapo{26--28} \support{Fiche · Act. 4} \support{Cahier}}}
{\pt{Réalise l’expérience de la cloche à vide (ou la projette, diapo 26) ; éclaire les deux plaques et fait relever les températures (diapo 27).}
 \pt{Fait traiter les questions 1 à 3 ; corrige au clic (diapo 28).}}
{\pt{Nomme le rayonnement ; compare : $+\dC{24,3}$ (noire) contre $+\dC{7,3}$ (blanche) ; distingue conduction, convection et rayonnement.}}
{\form Correction collective.\par
 \textit{Indicateur :} le rayonnement se propage dans le vide.}
{20}{1 h 24}
\etape{cAmber}{\axe{cRed}{ic_retenir}{À retenir}{Le rayonnement ; les trois modes de transfert}{\diapo{29} \support{Cahier}}}
{\pt{Projette la diapo 29 (conduction, convection, rayonnement animés) ; laisse le temps de recopier.}}
{\pt{Recopie l’« À retenir » dans son cahier.}}
{\form Question flash : « pourquoi les absorbeurs des chauffe-eau solaires sont-ils noirs ? »}
{8}{1 h 32}
\etape{cAmber}{\axe{cPurple}{ic_application}{Application}{La bouteille isolante}{\diapo{30} \support{Cahier}}}
{\pt{Montre une bouteille isolante ; projette l’énoncé (diapo 30) ; réponses dans le cahier ; correction au clic.}}
{\pt{Explique : le vide supprime conduction et convection ; les parois argentées réfléchissent le rayonnement.}}
{\somm Réponses rédigées.\par
 \textit{Indicateur :} un rôle pour chaque élément.}
{6}{1 h 38}
% ---------------------------------------------------------------------
\partie{cOrange}{Évaluation et bilan}{22 min · objectifs O2 à O6}
\etape{cGreen}{\axe{cOrange}{ic_exercices}{Évaluation}{Vrai ou faux ?}{\diapo{31}}}
{\pt{Projette le vrai/faux interactif (diapo 31) ; fait voter à main levée, puis justifier.}}
{\pt{Répond à main levée et corrige les affirmations fausses dans le cahier : a) F, b) V, c) F, d) V, e) V, f) F.}}
{\form Vote à main levée.\par
 \textit{Indicateur :} au moins 5 bonnes réponses sur 6.\par\vspace{.8mm}
 \regul Reprendre les affirmations les plus échouées.}
{6}{1 h 44}
\etape{cOrange}{\axe{cOrange}{ic_exercices}{Exercice 2}{Un morceau de glace dans un calorimètre}{\diapo{34--35} \support{Cahier}}}
{\pt{Projette l’énoncé (diapo 34) ; travail individuel dans le cahier ; corrige au clic (diapo 35).}}
{\pt{Calcule $Q_n = \qty{21294}{J}$ et $|Q_d| = \qty{53220}{J}$ ; conclut que toute la glace fond ; trouve $\theta_e \approx \dC{15,8}$.}}
{\somm Exercice individuel (évaluation de fin de chapitre).\par
 \textit{Indicateurs :} vérification préalable de la fusion complète ; trois termes pour la glace ; résultat avec unité.}
{12}{1 h 56}
\etape{cGreen}{\axe{cGreen}{ic_reponses}{L’essentiel}{Synthèse du chapitre}{\diapo{38}}}
{\pt{Projette « L’essentiel du chapitre » (diapo 38) ; donne le travail à la maison (encadré ci-dessous).}}
{\pt{Relit la synthèse ; note le travail à faire.}}
{\somm Bilan de l’évaluation.\par
 \textit{Exploitation :} repérer les notions à reprendre.}
{4}{2 h 00}
\end{xltabular}

\vspace{2mm}
\noindent
\begin{minipage}[t]{0.49\linewidth}\vspace{0pt}
\encadre{fi_maison}{cOrange}{Travail à la maison et prolongements}{%
\pt{\textbf{Après la séance 1} — exercice 1 (énoncé projeté, diapo 32, à recopier et à résoudre dans le cahier) : objectifs O4 et O6.}
\pt{\textbf{Après la séance 2} — exercice 3 (four solaire, énoncé projeté, diapo 36, à recopier et à résoudre dans le cahier) : objectifs O2, O5 et O7.}
\pt{Soutien : refaire l’application du chauffe-eau pour \qty{50}{L} chauffés de \dC{20} à \dC{60}. Approfondissement : dans l’exercice 2, calculer la masse maximale de glace à \dC{-10} que l’on peut faire fondre entièrement.}}
\end{minipage}\hfill
\begin{minipage}[t]{0.49\linewidth}\vspace{0pt}
\encadre{ic_application}{cPurple}{Bilan de la séquence (à compléter par le professeur)}{%
{\sffamily\footnotesize Objectifs atteints :}\lignevide\lignevide
\par\vspace{1.5mm}{\sffamily\footnotesize Difficultés rencontrées :}\lignevide\lignevide
\par\vspace{1.5mm}{\sffamily\footnotesize Ajustements pour la prochaine séquence :}\lignevide}
\end{minipage}

\end{document}
